% docs/manual/common/preamble.tex
% 由顶层 main_*.tex 在 \documentclass 之后 \input
% 与 styles.tex / macros.tex 平级，但只负责宏包导入与全局配置

% --- 链接与 PDF 书签 ---
\usepackage[
  bookmarks=true,
  bookmarksnumbered=true,
  bookmarksopen=true,
  colorlinks=true,
  linkcolor=blue!60!black,
  citecolor=blue!60!black,
  urlcolor=blue!60!black,
  pdfencoding=auto
]{hyperref}

% --- 紧凑页面布局（节约页数）---
\usepackage[a4paper, top=2.2cm, bottom=2.2cm, left=2cm, right=2cm, headheight=14pt]{geometry}

% --- 紧凑章节首页（titlesec 与 ctex 兼容）---
\usepackage[compact]{titlesec}
\titleformat{\chapter}[hang]
  {\Huge\bfseries}      % 字号样式
  {\thechapter\quad}    % 编号
  {0pt}                 % 编号与标题间距
  {\Huge}               % 标题前置
\titlespacing*{\chapter}{0pt}{-20pt}{16pt}
\titlespacing*{\section}{0pt}{14pt}{6pt}
\titlespacing*{\subsection}{0pt}{10pt}{4pt}

% --- 文本模式下让下划线 _ 不再触发 math subscript 报错 ---
% 写 \yamlkey{strict_reference} 这类带 _ 的字段时不必手动 \_ 转义
\usepackage{underscore}

% --- 表格 ---
\usepackage{booktabs}
\usepackage{longtable}
\usepackage{array}
\usepackage{tabularx}

% --- 图形 ---
\usepackage{graphicx}
\usepackage{tikz}
\usetikzlibrary{positioning,arrows.meta,shapes.geometric,backgrounds,fit,calc}

% --- 颜色 ---
\usepackage{xcolor}
\definecolor{obblue}{HTML}{0B5394}
\definecolor{obamber}{HTML}{C18B00}
\definecolor{obred}{HTML}{B23A3A}
\definecolor{obgreen}{HTML}{2E7D32}
\definecolor{obgray}{HTML}{4A4A4A}

% --- 等宽字体（minted 默认 lmmono 没有 ✓✗ 等 Unicode glyph，换 Menlo/DejaVu 兜底）---
\usepackage{fontspec}
\IfFontExistsTF{Menlo}{%
  \setmonofont{Menlo}[Scale=0.85]
}{%
  \IfFontExistsTF{DejaVu Sans Mono}{%
    \setmonofont{DejaVu Sans Mono}[Scale=0.85]
  }{}%
}

% --- 代码高亮 ---
\usepackage{minted}
\setminted{
  fontsize=\small,
  breaklines=true,
  breakanywhere=true,
  frame=lines,
  framesep=2mm,
  bgcolor=gray!5,
  numbersep=5pt,
  tabsize=2,
  autogobble=true,
}

% --- 标注框 ---
\usepackage[most]{tcolorbox}

% --- 页眉页脚 ---
\usepackage{fancyhdr}
\setlength{\headheight}{14pt}
\pagestyle{fancy}
\fancyhf{}
\fancyhead[LE,RO]{\thepage}
\fancyhead[RE]{\nouppercase\rightmark}
\fancyhead[LO]{\nouppercase\leftmark}
\renewcommand{\headrulewidth}{0.4pt}

% --- subfiles（顶层 manual.tex 用，subfile 内不会重复加载） ---
\usepackage{subfiles}

% --- 元数据（PDF 属性，被各卷覆盖） ---
\hypersetup{
  pdftitle={OpenBench 手册},
  pdfauthor={OpenBench 项目组},
  pdfsubject={陆面模型评估框架},
  pdfkeywords={OpenBench, land surface model, benchmarking, 中文}
}
